Sehr geehrter Herr Baier,

die Entwicklung hochpräziser medizintechnischer Systeme erfordert eine
strukturierte Brücke zwischen physikalischen Prinzipien und robuster,
serientauglicher Konstruktion. Als Physikingenieur mit Schwerpunkt auf
optomechatronischen Systemen bringe ich genau diese Schnittstellenkompetenz
mit. Während meiner Promotion an der Medizinischen Universität Innsbruck und
meiner Forschungsarbeit in Jena habe ich komplexe medizinische und optische
Geräte von der ersten Konzeption bis zur praktischen Anwendung im klinischen
Umfeld realisiert.

Ein zentraler Schwerpunkt meiner Arbeit liegt in der systematischen
Geräteentwicklung nach dem V Modell. So habe ich ein hochkomplexes
Multiphotonenmikroskop von Grund auf konzipiert, mechanisch und optisch mit
SolidWorks konstruiert und erfolgreich in Betrieb genommen. Dabei war die
Erstellung lückenloser Fertigungsunterlagen und technischer Dokumentationen
ebenso selbstverständlich wie die Durchführung von Risikoanalysen mittels
FMEA, um Fehlerquellen bereits in der Prototypenphase systematisch zu
minimieren.

Neben der reinen Konstruktion bringe ich fundierte Erfahrung in der
Qualitätssicherung und Messtechnik ein. Bei der Entwicklung eines faseroptischen
Raman Spektrometers für die Gewebecharakterisierung habe ich strukturierte
Testverfahren nach ISO 9001 und ISO 10110 etabliert. Durch die Kombination von
präziser Hardware Integration und intelligenter Datenverarbeitung in Python
konnte ich Systeme schaffen, die nicht nur im Labor, sondern auch in der
praktischen Anwendung verlässliche Ergebnisse liefern.

Ich schätze die Arbeit in interdisziplinären Teams und freue mich darauf, meine
Erfahrung in der Systementwicklung und Dokumentation bei W\&H einzubringen.
Meine Promotion werde ich im Juni 2026 abschließen. Mein österreichischer
Aufenthaltstitel wandelt sich danach nahtlos in eine Rot Weiss Rot Karte plus
um, sodass einem direkten und unkomplizierten Arbeitsbeginn in Bürmoos nichts
im Wege steht.
